\section{Patient Coding and the Similarity Score}
\label{app:coding}

This appendix describes how the six-attribute code of a patient is built, how predicted codes are generated, and how the similarity score is computed, with a worked example. The scheme follows the classification-and-coding approach of group technology, in which an item is described by a short string of coded attributes and items are compared by a similarity coefficient \cite{OFFODILE1990,Seifoddini01091986,YIN2005471}.

\subsection{Attributes and Classes}

Table~\ref{tab:coding} lists the attributes. Attribute A is the CTAS level assigned at triage and is observed without error. Attributes B to F describe the patient's expected course in the ED; in practice they would come from a prediction model at triage, and in the simulation they are predicted from the patient's true attributes as described below.

\begin{table}[!htbp]
\centering
\caption{Coding scheme. Lower codes indicate higher acuity, longer service, admission, diagnostic need, heavier resource use or consultation.}
\label{tab:coding}
\footnotesize
\begin{adjustbox}{max width=\linewidth}\begin{tabular}{lp{3.6cm}p{6.2cm}}
\toprule
Attribute & Classes & Definition in the model \\
\midrule
A: CTAS level & 1--5 & Triage level drawn from the acuity mix (observed). \\
B: Service class & 1 very long, 2 long, 3 average, 4 short, 5 very short & Quintile of total treatment time (physician assessment and reassessment, diagnostic and consultation steps, discharge process; waits and boarding excluded). Cut points \SvcCutOne, \SvcCutTwo, \SvcCutThree{} and \SvcCutFour{} minutes, from the community ED. \\
C: Admission & 1 admitted, 2 discharged & Disposition drawn from the CTAS-specific admission probability. \\
D: Diagnostics & 1 yes, 2 no & At least one laboratory, X-ray or CT step in the pathway. \\
E: Resource use & 1 three or more, 2 two, 3 one, 4 none & Number of diagnostic and consultation steps in the pathway. \\
F: Consultation & 1 yes, 2 no & Consultation (DR) step in the pathway. \\
\bottomrule
\end{tabular}\end{adjustbox}
\end{table}

\subsection{Predicted Codes}

For each predicted attribute with $n$ classes, a reference confusion matrix $M_{\text{ref}}$ has the reference accuracy on its diagonal; the remaining probability of row $i$ is spread over the other classes $j$ in proportion to $\delta^{|i-j|-1}$ with $\delta=\ErrDecay$, so that errors fall mostly on neighbouring classes. Table~\ref{tab:confusion} shows the two ordinal matrices. The matrix used at error level $\alpha$ is $M(\alpha)=I+\alpha(M_{\text{ref}}-I)$, which is a valid stochastic matrix for every $\alpha$ for which its diagonal is non-negative. A patient's predicted class is obtained by inverse-transform sampling from row $p$ of $M(\alpha)$ with one uniform random number per patient and attribute; because the same number is used at every $\alpha$, the predicted class changes monotonically from the true class as $\alpha$ grows, and comparisons across $\alpha$ use common random numbers.

\begin{table}[!htbp]
\centering
\caption{Reference confusion matrices ($\alpha=1$) for the two ordinal predicted attributes. Rows: true class; columns: predicted class.}
\label{tab:confusion}
\footnotesize
\begin{adjustbox}{max width=\linewidth}\begin{tabular}{lrrrrr}
\toprule
True class & \multicolumn{5}{c}{Predicted class (service time; 1 very long \dots 5 very short)} \\
\cmidrule(lr){2-6}
 & 1 & 2 & 3 & 4 & 5 \\
\midrule
1 & 0.450 & 0.363 & 0.127 & 0.044 & 0.016 \\
2 & 0.222 & 0.450 & 0.222 & 0.078 & 0.027 \\
3 & 0.071 & 0.204 & 0.450 & 0.204 & 0.071 \\
4 & 0.027 & 0.078 & 0.222 & 0.450 & 0.222 \\
5 & 0.016 & 0.044 & 0.127 & 0.363 & 0.450 \\
\midrule
 & \multicolumn{5}{c}{Predicted class (resource use; 1 three or more steps \dots 4 none)} \\
\cmidrule(lr){2-5}
 & 1 & 2 & 3 & 4 & \\
\midrule
1 & 0.550 & 0.306 & 0.107 & 0.037 & \\
2 & 0.191 & 0.550 & 0.191 & 0.067 & \\
3 & 0.067 & 0.191 & 0.550 & 0.191 & \\
4 & 0.037 & 0.107 & 0.306 & 0.550 & \\
\bottomrule
\end{tabular}\end{adjustbox}
\end{table}

\subsection{Similarity Score and Queue Order}

For a code $p$, a target code $q$ and weights $w\geq 0$, the score is (\mref{eq:sim}): $S(p)=\sum_k w_k(1-|p_k-q_k|/R_k)$ with $R_k=n_k-1$. Each term lies between $0$ and $w_k$, so with weights that sum to one the score lies in $[0,1]$. A queue ordered by similarity serves CTAS~1 patients first and then the patient with the largest score; the score is rounded to $10^{-9}$ and ties are broken by the time the patient joined the queue and then by patient identifier. With an ageing term $\gamma$, the effective priority at time $t$ of a CTAS-$l$ patient who joined at $t_0$ is $S(p)+\gamma(t-t_0)/T_l$, where $T_l$ is the CTAS target. With a reassessment bonus $\beta$, a patient waiting for a physician reassessment receives $S(p)+\beta$ in the physician queue.

\subsection{Worked Example}

Table~\ref{tab:example} scores four patients. Under the SIM rule (target $1$-$5$-$2$-$2$-$4$-$2$, weights \SimWA, \SimWB, \SimWC, \SimWD, \SimWE, \SimWF), patient P3, a CTAS~3 patient expected to be discharged quickly without tests, has score \WexPThreeSim{} and is ranked \WexPThreeRankSim{}, ahead of the CTAS~2 patients P1 and P2. Under the exact representation of CTAS-SPT ($w_A=\RepSptWA$, $w_B=\RepSptWB$), both CTAS~2 patients are ranked first, and among them the patient with the shorter expected service (P2) comes first, as Proposition~\mref{prop:lex} requires. The last two columns show similarity routing: the difference $S_{\text{second}}-S_{\text{main}}$ is positive for P3 and P4, which go to the second area; P1 and P2 are CTAS~2 and stay in the Main area by the safety rule.

\begin{table}[!htbp]
\centering
\caption{Worked example: scores and ranks of four patients under SIM and under the exact similarity-index representation of CTAS-SPT, and similarity routing (values computed with the model code).}
\label{tab:example}
\footnotesize
\begin{adjustbox}{max width=\linewidth}\begin{tabular}{lcrrrrrl}
\toprule
Patient & Code A-B-C-D-E-F & $S_{\text{SIM}}$ & Rank & $S_{\text{rep}}$ & Rank & $S_{\text{second}}-S_{\text{main}}$ & Area \\
\midrule
P1 & 2-1-1-1-1-1 & 0.3000 & 4 & 0.6375 & 2 & -0.8000 & Main \\
P2 & 2-4-2-1-3-2 & 0.7083 & 2 & 0.7500 & 1 & +0.0167 & Main \\
P3 & 3-5-2-2-4-2 & 0.8000 & 1 & 0.5750 & 3 & +0.6000 & Second \\
P4 & 4-3-2-1-2-2 & 0.4167 & 3 & 0.2875 & 4 & +0.2333 & Second \\
\bottomrule
\end{tabular}\end{adjustbox}
\end{table}

\subsection{Why the Weights Need to Be Checked}

The worked example shows how a weighting that is read as ``acuity first'' can rank a less urgent patient first. With weights $(\SimWA,\SimWB,\dots)$ the largest possible gain from one CTAS level is $\SimWA/4$, while the other attributes together can add up to the sum of their weights. Proposition~\mref{prop:lex} turns this observation into a condition that can be checked before a priority index is deployed, and Section~\mref{sec:optimization} shows how weights can instead be chosen by simulation-based optimization.
